\section{Construct a Binary Tree from Postorder and Inorder}
\label{sec:trees-construct-post-in}

\problemstatement{Given the postorder and inorder traversals of a binary
tree with unique values, reconstruct and return the tree.}

\begin{patternbox}
This is the previous section mirrored. Postorder visits left, right,
\emph{root} -- so the root is the \textbf{last} element, not the first. Read
postorder \emph{right to left} and it behaves like preorder with the
children swapped: root, then right subtree, then left subtree. So build the
\textbf{right subtree before the left}, consuming postorder from the back.
\end{patternbox}

\subsection*{Understanding the problem carefully}

Because postorder ends with the root, a cursor moving backward through
postorder yields roots in the order root, right-subtree-root,
\ldots. Inorder still splits into left and right around the root's
position. The one change from the preorder version: since we now encounter
the right subtree's root immediately after the current root (reading
backward), we must recurse \emph{right first}, then left, so the backward
cursor lines up with the subtrees correctly.

\begin{ideabox}[Postorder from the Back Names Roots; Build Right, Then Left]
Take the current root as postorder[\textit{postIdx}] and decrement
\textit{postIdx}. Split inorder at the root's position. Recurse into the
\emph{right} subtree first (node.right), then the \emph{left}
(node.left) -- the reverse of the preorder construction.
\end{ideabox}

\subsection*{Algorithm}

\begin{enumerate}
  \item Build \textit{pos}: value $\to$ inorder index. Set
        \textit{postIdx} $=n-1$.
  \item \textsc{Build}(\textit{inLo}, \textit{inHi}):
  \begin{enumerate}
    \item If \textit{inLo} $>$ \textit{inHi}, return null.
    \item \textit{rootVal} $=$ postorder[\textit{postIdx}--]; make node.
    \item $p=\textit{pos}[\textit{rootVal}]$.
    \item node.right $=$ \textsc{Build}($p+1$, \textit{inHi}); node.left
          $=$ \textsc{Build}(\textit{inLo}, $p-1$).
    \item Return node.
  \end{enumerate}
\end{enumerate}

\subsection*{Dry run}

\dryrun{postorder $=[9,15,7,20,3]$, inorder $=[9,3,15,20,7]$.}

\begin{itemize}
  \item Root $=3$ (postorder last); inorder index $1$: left $=[9]$, right
        $=[15,20,7]$.
  \item Build right first: next root (moving back) $=20$; inorder index
        $3$: right $=[7]$ (leaf), left $=[15]$ (leaf).
  \item Build left: next root $=9$ -- leaf.
  \item Reconstructed tree matches the original.
\end{itemize}

\subsection*{C++ implementation}

\begin{lstlisting}
#include <bits/stdc++.h>
using namespace std;

struct TreeNode {
    int val;
    TreeNode* left;
    TreeNode* right;
    TreeNode(int v) : val(v), left(nullptr), right(nullptr) {}
};

int postIdx;

TreeNode* build(vector<int>& postorder, int inLo, int inHi,
                unordered_map<int,int>& pos) {
    if (inLo > inHi) return nullptr;
    int rootVal = postorder[postIdx--];
    TreeNode* node = new TreeNode(rootVal);
    int p = pos[rootVal];
    node->right = build(postorder, p + 1, inHi, pos);  // right first
    node->left  = build(postorder, inLo, p - 1, pos);
    return node;
}

TreeNode* buildTree(vector<int>& inorder, vector<int>& postorder) {
    unordered_map<int,int> pos;
    for (int i = 0; i < (int)inorder.size(); i++) pos[inorder[i]] = i;
    postIdx = postorder.size() - 1;
    return build(postorder, 0, inorder.size() - 1, pos);
}

void inorderPrint(TreeNode* n) {
    if (!n) return;
    inorderPrint(n->left);
    cout << n->val << " ";
    inorderPrint(n->right);
}

int main() {
    vector<int> inorder   = {9, 3, 15, 20, 7};
    vector<int> postorder = {9, 15, 7, 20, 3};
    TreeNode* root = buildTree(inorder, postorder);
    inorderPrint(root);
    cout << "\n";
    return 0;
}
\end{lstlisting}

\begin{complexitybox}
\textbf{Time Complexity.} $O(n)$ with the inorder index map. \textbf{Space
Complexity.} $O(n)$ for the map plus $O(h)$ recursion.
\end{complexitybox}

\begin{takeawaybox}
Postorder+inorder is preorder+inorder read backward with the child order
flipped. Seeing the symmetry means you derive one from the other instead of
memorising two separate algorithms. Note: preorder+postorder alone cannot
uniquely rebuild a tree -- inorder is what fixes the left/right split.
\end{takeawaybox}
